\documentclass[a4paper, 11pt]{article}

\usepackage[utf8]{inputenc}
\usepackage{amsmath}
\usepackage{amsfonts}
\usepackage{amssymb}
\usepackage{graphicx}
\usepackage{amsthm}
\usepackage{enumerate}
\usepackage{verbatim}
\usepackage{hyperref}
%\usepackage{tkz-graph}
%\usepackage{algorithm}
%\usepackage{algorithmic}
\usepackage{array,epsfig}
\usepackage{amsxtra}
\usepackage{mathrsfs}
\usepackage{color}
\usepackage{float}
%\usepackage{subfigure}
\usepackage{caption}
\usepackage{subcaption}
%\usepackage{wrapfig}
\usepackage[a4paper]{geometry}



\begin{document}

\begin{center}
{\LARGE{Project 1 \\[12pt] Name1, Name2, and Name3}}
\end{center}


\section*{Problem 1}

\subsection*{a)}
This is an equation
\[
    \mathbf{P} = \begin{bmatrix} 0 & 0 & 1 \\
                                 1 & 0 & 0 \\
                                 0 & 1 & 0 \\
                 \end{bmatrix}
\]
a numbered equation
\begin{equation}
    \pi_0 = 0.1 \label{eq:pi0},
\end{equation}
where Equation \eqref{eq:pi0} shows that...

\subsection*{b)}


\begin{figure}[ht]
 \begin{center}
% \includegraphics[width=12cm]{res1b_1.png}
 \end{center}
 \caption{Plot of realizations.} \label{fig1}
 \end{figure}
Figure \ref{fig1} shows 10 realizations.



\section*{Problem 2}

\subsection*{a)}
This is the answer to 2a).\\

\noindent
This is another paragraph


\end{document}


